\section{Proof of the interpolation theorem}\label{int:proof}

We prove \Cref{int:height} by the graph construction explained in
\Cref{int:section}.  The algebraic bounds below control every new
variety and every ideal-membership certificate.  An ordinary differential
estimate supplies the pole along each normalized boundary divisor;
layered interpolation on its graph then gives the polynomial used in
the product-formula descent.

\subsection{Effective algebra with absolute heights}

We use polynomial bounds in a fixed number of variables.  More precisely,
in the lemmas below each assertion of such a bound means an inequality
$c(1+D)^c$ or $c(1+D+T)^c(1+H)$; the constant and exponent depend only
on the fixed number of variables.  After a fixed number of operations these
are still polynomial bounds.  Heights are logarithmic throughout.

\begin{lemma}[Algebraic degree and height bounds]\label{int:algebra}
Fix $M$.  Suppose at most $D$ polynomials in $M$ variables have degree at
most $D\geq2$ and affine coefficient height at most $H$.
\begin{enumerate}[label=\textup{(\alph*)}]
\item Their ideal has a degree-compatible Gr\"obner basis of polynomially
bounded degree and coefficient height at most
$\operatorname{poly}_M(D)(1+H)$.  Each basis element has a representation
in the original generators whose degree and coefficient height satisfy
the same types of bounds.  Normal-form reduction can be chosen linear.
For an input of degree $T$, its representation by the remainder and the
original generators has degree at most $T+\Delta$, where
$\Delta=\operatorname{poly}_M(D)$ is independent of $T$.
\item If a polynomial of degree $T$ belongs to the ideal, it has a
certificate of membership of degree polynomial in $D+T$.  The coefficient
height of a certificate can be bounded polynomially in $D+T$ times one
plus the heights of the given polynomial and generators.
\item Every irreducible component of the common zero set has degree
polynomial in $D$, and admits set-theoretic defining equations whose number
and degrees are polynomial in $D$ and whose affine coefficient heights
are at most $\operatorname{poly}_M(D)(1+H)$.
\end{enumerate}
All bounds are absolute and are unchanged by enlarging the number field.
The ideals in \textup{(a)} and \textup{(b)} need not be radical.
\end{lemma}

\begin{proof}
The fixed-variable degree inputs are the affine Gr\"obner basis degree
bound of \cite[Corollary~8.3]{Dube1990} and the polynomial
ideal-membership degree bound of
\cite[Theorem~3.4]{Aschenbrenner2004}, which gives a proof of the
Hermann--Seidenberg method.  Their degree bounds become
polynomials in $D$ when $M$ is fixed.  For example, using one additional
homogenizing variable, a larger permissible bound for the first input is
\[
 \Gamma_M(D)=\left\lceil2(D^2/2+D)^{2^M}\right\rceil.
\]
Only this polynomial character is needed here.

We spell out the height consequence, since it must not introduce the
degree or discriminant of a field of definition.  A polynomial identity
of a prescribed degree becomes a linear system by equating coefficients.
There are at most $\binom{T+M}{M}$ monomials of degree at most $T$.
After the degree bounds just quoted, the number of entries of every such
system is polynomial in $D+T$.  For a matrix with $u$ entries, each of
height at most $H_1$, the affine height of the full matrix is at most
$uH_1$.  A determinant of size $r$ then has height at most
$r uH_1+\log(r!)$: this follows place by place from the determinant
expansion, using the ultrametric inequality at finite places.  Choosing
a nonzero pivot minor and applying Cramer's rule gives solution and kernel
bases of height polynomial in the matrix size times $1+H_1$.
The product formula for the pivot denominator proves these assertions
with normalized absolute heights.  There is no factor $[F:\Q]$.

The reduced Gr\"obner basis can be recovered by these linear systems:
each required polynomial has bounded degree, and membership in the
generator ideal can be imposed with the bounded-degree certificates.
Leading coefficients and the coefficients of forbidden leading monomials
are fixed by linear equations.  This gives the stated height bounds for
the basis and its representations.  Fix those representations.  With a
degree-compatible order, division of a monomial of degree at most $T$
uses multiples of basis elements of degree at most $T$.  Replacing a
basis element by its fixed representation increases this degree by at
most a fixed additive amount $\Delta$.  Reducing each monomial and
extending linearly gives the required linear reduction map.  This proves
\textup{(a)}; the same coefficient argument and the membership degree
bound prove \textup{(b)}.

For \textup{(c)}, pass to projective closures.  Follow a component of
interest by proper hypersurface cuts, starting with projective space.
Whenever the current component is larger than the desired one, some
given defining polynomial is nonzero on the current component; cut by
it and choose a component still containing the desired one.  There are
at most $M$ proper cuts.  Geometric B\'ezout bounds the degrees at each
step.  For heights, apply the global intersection bound
\cite[Corollary~2.11]{KrickPardoSombra2001} to the affine parts of these
cuts.  It bounds the height of each affine intersection by a polynomial
in the degrees times one plus the preceding height and the coefficient
heights of the cutting polynomials.  The height of an affine variety
in that result is defined through its projective closure
\cite[\S1.2.4]{KrickPardoSombra2001}.  This absolute Chow height is
nonnegative and additive on components, so it bounds each retained
component separately.  Its comparison with the projective coefficient
height of a Chow form costs a polynomial
in the degree; this follows also by comparing the coefficient norm with
the logarithmic integral norm in the definition of Chow height.
Components supported at infinity do not enter these affine intersections.

Here is an explicit way to return to equations.  Let $X$ be the projective
closure of a component, of dimension $r$ and degree $\delta$, and let
$\Phi_X$ be its Chow form.  In each of its $r+1$ hyperplane groups
substitute the coefficients of a hyperplane through $[1:x]$:
\[
 (u_{i0},u_{i1},\ldots,u_{iM})
   =\left(-\sum_{j=1}^M u_{ij}x_j,u_{i1},\ldots,u_{iM}\right).
\]
Set every coefficient in the remaining $u_{ij}$ equal to zero.  These
are polynomials in $x$, of degree at most $(r+1)\delta$, with
polynomially many coefficients.  They all vanish if $[1:x]\in X$.
If $[1:x]\notin X$, projection from this point has image of dimension
at most $r$; $r+1$ general hyperplanes through the point have no common
intersection with $X$.  The substituted Chow form is then not the zero
polynomial.  Thus the coefficient equations define exactly the affine
component.  Normalize a nonzero coefficient of $\Phi_X$ to one.
Its affine coefficient height is its projective coefficient height by
the product formula, and the indicated substitutions increase height
by a polynomial in $\delta$.  This proves the equation-height and
equation-count assertions, including after a field extension needed to
define the component.
\end{proof}

\subsection{A uniform boundary multiplicity bound}

The next lemma supplies a truncation $F=E_{<m}$ for which $F/a^m$ has
a pole along every normalized boundary divisor.  This is the input to
\Cref{int:graph}; an identity or nonidentity merely on the divisor
would not determine that pole.  Only degrees enter the estimate.  In
particular, the coefficients of the transverse slices used in its proof
need no height bound.
The differential method uses the span of successive derivatives;
compare \cite[Theorem~1]{BertrandBeukers1985}. We prove the version
needed here because the algebraic curve carrying the pulled-back
system varies during the dimension descent.

\begin{lemma}[Ordinary differential multiplicity]\label{int:multiplicity}
Use the fixed function data of \Cref{int:height}.  There is a polynomial
$M_0(D)$, depending only on those data and $N$, with the following property.
Let $W\subseteq\mathbb A^N$ be irreducible of degree at most $D$, with
$a$ nonconstant, and suppose its normalized boundary branches satisfy
condition \textup{(iv)}.  There are integers
$0\leq c_\alpha\leq M_0(D)$ and an integer
$1\leq m\leq M_0(D)$ such that, for
\[
 E=\sum_{\alpha=1}^s c_\alpha E_\alpha
   =\sum_{j\geq0}a^j e_j(U),
 \qquad E_{<m}=\sum_{j=0}^{m-1}a^je_j(U),
\]
every prime divisor $D_0$ of the normalization over $W\cap\{a=0\}$
satisfies
\begin{equation}\label{int:strict-order}
 \ord_{D_0}E_{<m}<m\ord_{D_0}a.
\end{equation}
The constants do not depend on the coefficient sizes of $W$ or of the
polynomial combinations defining $E_\alpha$.
\end{lemma}

\begin{proof}
Put $k=\dim W$ and let $\nu:\widetilde W\to W$ be its finite
normalization.  Since $a$ is a nonzero nonunit whenever the boundary is
nonempty, its zero set is pure of codimension one.  Finiteness of $\nu$
implies that every divisor of $a$ upstairs maps onto a component of this
zero set.  At a general point of such a divisor $D_0$, both
$\widetilde W$ and $D_0$ are smooth, and the map on $D_0$ has rank $k-1$.
Indeed the induced extension of function fields is finite and separable
in characteristic zero.

Choose $k-1$ general affine hyperplanes in the original ambient space.
Their pullbacks cut smooth curve germs transverse to $D_0$ at its general
intersection points.  They can also be chosen to avoid the zero set of
the first nonzero transverse coefficient of any nonzero $E_\alpha$.
To see the latter assertion explicitly, use local coordinates $(t,s)$
with $D_0=\{t=0\}$.  A nonzero germ has the expansion
$t^r u(s)+O(t^{r+1})$ with $u$ nonzero on an open subset of $D_0$.
A transverse curve through a point of that open subset preserves its
order $r$.  Infinite order would instead mean that all transverse
coefficients vanish on a local open set; analyticity would make the
germ identically zero.  Thus the hypothesis supplies a finite first
order on each branch to be tested.  The excluded set may be taken
Zariski closed: once this finite order is known, choose a truncation
longer than its ratio to $\ord_{D_0}a$.  That algebraic truncation
has the same leading transverse coefficient.  Its nonvanishing is
the nonvanishing of an algebraic section of a power of the conormal
line bundle on a smooth open subset of $D_0$.  This justifies using
general algebraic hyperplanes to make all these choices simultaneously.

The sliced curve has degree at most $D$, after removing components that
do not meet the chosen open sets.  The bound on the number of branches
can be seen directly on the finite normalization $\overline\nu$ of the
projective closure of $W$.  Put $L=\overline\nu^*\mathcal O(1)$.
The homogeneous coordinate defining $a=0$ is a nonzero section of $L$;
if its divisor is $\sum_T m_TT$, then
\[
 \deg W=L^k
   =\sum_T m_T\,[\Qbar(T):\Qbar(\overline\nu(T))]
                     \deg\overline\nu(T).
\]
Each term is a positive integer.  Every affine normalized divisor
under consideration occurs in this sum, so there are at most $D$ of
them.  Equivalently the transverse curve branches are counted with
their positive hyperplane-intersection multiplicities.  This also
explains why no bound for equations of $\widetilde W$ is needed.

On a tested curve, adjoin the algebraic curve coordinates implicit in
the fixed germs.  A parameter on each fixed curve gives a fixed finite
map of its smooth projective completion to $\mathbb P^1$.  Pulling
these maps back by the functions $a$ or $aR$ and taking the component
of the specified germ increases degree by a fixed factor.  Repeating
this for the fixed finite list still gives an algebraic curve of degree
polynomial in $D$ in a fixed product of projective spaces.  A fixed
Segre embedding and normalization produce a smooth projective curve
$X$ with genus polynomial in $D$.  For example a general plane
projection bounds its genus by the arithmetic genus of a plane curve
of polynomial degree.  The degrees of $a$, $R$, the original coordinate
functions, and each adjoined coordinate on $X$ are polynomial in $D$.
The chosen points of $X$ lie over the specified ordinary centers.
Ramification in these maps does not destroy regularity of the pulled
back connections.  If an argument is constant, its germ contributes
a constant solution vector instead.

Direct sums, tensor products, and duals encode every degree-$b$
polynomial combination of the fixed germs as
\begin{equation}\label{int:row-solution}
 f=r\,y,
\end{equation}
where $r$ is a rational row on $X$ and $y$ is a horizontal analytic
vector for one algebraic connection of rank at most a fixed integer
$r_*$.  The coordinates and the target-dependent scalar coefficients
occur in $r$.  The integer $r_*$ is determined before $W,d,h$ are chosen:
it is the sum of the ranks of the finitely many tensor and dual systems
needed for all monomials of degree at most $b$.  Choose the tensor
rational frame obtained from the specified fixed frames.  It is regular
and invertible at the tested point, and the connection is ordinary
there.  The polar degrees of its rational connection matrix and of $r$
are polynomial in $D$, since all constituent rational maps and
functions have those degree bounds.

We give the zero estimate for \eqref{int:row-solution}.  Let
$\delta=d/da$ on the function field of $X$.  If $g$ is rational, then
\[
 (\delta g)_\infty
  \leq (g)_\infty+((g)_\infty)_{\mathrm{red}}+(da)_0,
 \qquad
 \deg(da)_0\leq 2g(X)+2\deg(a).
\]
The first inequality follows by writing $\delta g=dg/da$ in a local
parameter; the second uses the canonical divisor degree and the bound
$\deg(da)_\infty\leq2\deg(a)$.  Write the horizontal equation as
$\delta y=B y$ and define derivative rows
\[
 r_0=r,\qquad r_{j+1}=\delta r_j+r_jB.
\]
Then $r_jy=\delta^jf$.  The entries of $B$ have a common polar divisor
of polynomial degree: divide the rational connection one-form by
$da$, whose zero divisor has just been bounded.  If $P_j$ bounds the
poles of the entries of $r_j$ and $Q$ bounds the poles of $B$, then
\[
 P_{j+1}=2P_j+(da)_0+Q
\]
is a valid common bound.  For $j\leq r_*$ all these divisors therefore
have polynomial degree, independently of coefficient sizes.

Let $s'$ be the first index for which $r_{s'}$ lies in the function-field
span of $r_0,\ldots,r_{s'-1}$.  Unless $r=0$, we have $1\leq s'\leq r_*$.
Differentiating this dependence shows that the span is stable under
covariant differentiation; all subsequent rows lie in it.  Let $K_\eta$
be its common kernel.  At the tested point take the local ring of
convergent germs $\mathcal O=\C\{t\}$, the ordinary frame
$\mathcal O^{r_*'}$ of the actual rank $r_*'\leq r_*$, and the
saturation
\[
K=\bigl(K_\eta\otimes_{\Qbar(X)}\operatorname{Frac}\mathcal O\bigr)
       \cap\mathcal O^{r_*'}.
\]
This is a direct summand, since the quotient is torsion free over a
discrete valuation ring.  It is preserved by the regular connection
with respect to $d/dt$: generic preservation follows from that for
$\delta$, and the derivative of a regular section is regular.  Thus
the quotient has an ordinary connection.  If $f$ is nonzero, the
projection of $y$ to the quotient is nonzero and cannot vanish at
$t=0$, by uniqueness for an ordinary linear differential equation.
It is therefore primitive in the quotient lattice.

This primitivity can be used in the original frame, without a height
or pole bound on any quotient frame.  Let $A$ be the matrix with rows
$r_0,\ldots,r_{s'-1}$, and let $P$ be the degree of a common polar
divisor for its entries.  Its entries have order at least $-P$ at
$t=0$.  Some $s'\times s'$ minor $\Delta$ is nonzero, with polar
degree at most $s'P$, and hence with
$\ord_0\Delta\leq s'P$.  Smith normal form over $\mathcal O$, after
clearing the entry poles, has nonzero exponents
$\sigma_1\leq\cdots\leq\sigma_{s'}$ satisfying
\[
 \sigma_i\geq-P,\qquad
 \sum_{i=1}^{s'}\sigma_i
   =\min_{\Delta'\ne0}\ord_0\Delta'
   \leq s'P.
\]
Consequently $\sigma_{s'}\leq(2s'-1)P$.  The invertible row and column
matrices in Smith normal form preserve the local lattices and their
maximal ideals.  The saturated kernel corresponds to the zero columns,
so primitivity modulo $K$ means that some coordinate in a nonzero
column is a unit.  It follows that
\begin{equation}\label{int:smith-order}
 \min_{0\leq j<s'}\ord_0(\delta^j f)
   \leq (2s'-1)P.
\end{equation}

If $e=\ord_0 a\geq1$ and $n_0=\ord_0 f$, local differentiation gives
\[
 \ord_0(\delta^j f)\geq n_0-je.
\]
Cancellation of leading terms only increases the order, so the
inequality also covers that case.  Combining it with
\eqref{int:smith-order} bounds $n_0/e$ by a polynomial in $D$.
Finite pullback multiplies both orders by its ramification index;
the same bound holds on the original normalized divisor.  The
transverse slice was chosen to preserve its first nonzero coefficient,
so this proves the desired order bound for each nonzero original
expression on each divisor.

For each divisor take the least order among the nonzero $E_\alpha$.
Cancellation at that order in $\sum c_\alpha E_\alpha$ is a proper
linear condition on $(c_1,\ldots,c_s)$, with coefficients in a residue
function field.  There are at most $D$ such conditions.  A proper
hyperplane meets the grid $\{0,\ldots,D\}^s$ in at most
$(D+1)^{s-1}$ points: fix all coordinates except one whose coefficient
is nonzero.  The union of the forbidden hyperplanes cannot fill the
grid.  Choose a vector outside it.  The resulting $E$ has the stated
polynomial order bound on every divisor.  Choose an integer $m$ larger
than all the bounded order ratios.  All coordinate functions are regular
at these divisors, and
$E-E_{<m}=a^m\sum_{j\geq m}a^{j-m}e_j(U)$ is an analytic multiple
of $a^m$.  Its order is at least $m\ord a$.  The smaller order of
$E$ therefore equals that of $E_{<m}$, proving
\eqref{int:strict-order} and the lemma.
\end{proof}

\subsection{The graph and its projected boundary}

\begin{lemma}[Graph boundary]\label{int:graph}
Let $W$ be an irreducible affine variety of dimension $k\geq1$, and
let $a$ be a nonconstant regular coordinate.  Suppose a polynomial $F$
and an integer $m\geq1$ satisfy
\[
 \ord_D F<m\ord_Da
\]
on every normalized divisor above $W\cap\{a=0\}$.  Let $W'$ be the
affine closure of the graph
\[
 q=g:=a^{-m}F\quad\text{over }W\setminus\{a=0\},
\]
and let $\pi:W'\to W$ be the projection.  Then
\begin{equation}\label{int:projected-boundary}
 \dim\overline{\pi(W'\cap\{a=0\})}\leq k-2.
\end{equation}
In particular, for any polynomial $f(U)$,
\begin{equation}\label{int:graph-cut}
 \dim\bigl(W'\cap\{a=0, q+f(U)=0\}\bigr)\leq k-2.
\end{equation}
When $k=1$ these sets are empty.  Degrees and heights of $W'$ admit
the bounds of \Cref{int:algebra} in terms of defining equations of
$W$, $F$, and $m$.
\end{lemma}

\begin{proof}
The graph is irreducible of dimension $k$.  Over $a\ne0$ it is the
whole zero set of $a^mq-F$ in $W\times\mathbb A^1$, and its closure
is an irreducible component of that zero set.  This proves the
complexity assertion by \Cref{int:algebra}.

Let $\nu:\widetilde W\to W$ be the finite normalization and let
$\widetilde\Gamma$ be the closure of the same rational graph in
$\widetilde W\times\mathbb A^1$.  Its map to $W\times\mathbb A^1$
is finite, so its image is closed.  That image equals $W'$: it contains
the dense original graph and is contained in its closure.  At the
generic point of each normalized boundary divisor $D$, the reciprocal
$g^{-1}$ is regular and vanishes, because $g$ has a pole.  On a
neighborhood of that point, the upstairs graph closure satisfies
\[
 qg^{-1}=1.
\]
It therefore has no point above a dense open subset of $D$.

Fix a boundary component $B\subset W$.  All points of the finite
normalization above its generic point lie on the finitely many divisors
dominating $B$.  For each of these divisors remove the proper closed
subset not covered by the neighborhood just used.  The images of the
removed subsets are closed under the finite map and have dimension
at most $k-2$.  Outside their union, and outside intersections with
the other boundary components, no upstairs graph point lies over $B$.
The finite surjection $\widetilde\Gamma\to W'$ then gives the same
exclusion downstairs.  Doing this for every $B$ proves
\eqref{int:projected-boundary}.

On $q+f(U)=0$ the projection has inverse
$(a,U)\mapsto(a,U,-f(U))$.  Thus the set in
\eqref{int:graph-cut} is isomorphic to a closed subvariety of the
projected boundary and has its asserted dimension.  This step fixes
the graph coordinate; no bound on the dimension of
$W'\cap\{a=0\}$ itself was used.
\end{proof}

\subsection{The auxiliary polynomial in the actual ideal}

The following construction keeps a high power of $a$ in the membership
certificate.  Working with the actual defining ideal is essential for
retaining that power when the certificate is evaluated.

\begin{lemma}[Layered interpolation]\label{int:auxiliary}
Fix $M$.  In $\Qbar[a,U,q]$, with $M$ variables in total, let
$I=(I_1,\ldots,I_r)$ be generated by polynomials of controlled number,
degree, and height.  Suppose its zero set $W'$ is irreducible of
dimension $k\geq1$.  Let
\[
 f_0=q+e_m(U),\qquad J=(I,a,f_0),\qquad
 F_n=f_0+\sum_{j=1}^{n-1}a^j e_{m+j}(U),
\]
where $\deg e_{m+j}\leq C_0(1+m+j)$ and
$\dim V(J)\leq k-2$, interpreted as emptiness for $k=1$.
Let $D\geq2$ bound $r$, $m$, the generator degrees, and $C_0$.
There are constants $B,B'$, polynomial in $D$ for fixed $M$, such that
the following holds.  Whenever $n,L\geq1$ and
\begin{equation}\label{int:condition-count}
 B'n(1+L+Bn+B)^{k-2}<\binom{L+k}{k},
\end{equation}
there is a polynomial $P$ with $\deg P\leq L$, nonzero on $W'$, and
\begin{equation}\label{int:actual-membership}
 P\in(I,a^n,F_n).
\end{equation}
For $k=1$ the left side of \eqref{int:condition-count} is replaced by
zero.  Both $P$ and a certificate for \eqref{int:actual-membership}
can have degree polynomial in $D+L+n$ and coefficient height
polynomial in $D+L+n$ times one plus the coefficient heights of
$I_1,\ldots,I_r,e_m,\ldots,e_{m+n-1}$.

In particular, for any prescribed positive integer $c$, one can take
$n=cL$ with $L$ polynomial in $D+c$.
\end{lemma}

\begin{proof}
Fix a degree-compatible Gr\"obner basis of $J$ and bounded-degree
representations of its elements in $I_1,\ldots,I_r,a,f_0$ using
\Cref{int:algebra}.  Reducing each monomial and extending linearly
gives fixed linear maps, with
\begin{equation}\label{int:linear-division}
 D_1=\operatorname{NF}_J(D_1)
       +\sum_{\rho=1}^r I_\rho A_\rho(D_1)
       +aA(D_1)+f_0H(D_1).
\end{equation}
If $\deg D_1\leq T$, every coefficient polynomial on the right
has degree at most $T+\Delta$, where $\Delta$ is polynomial in $D$
and independent of $T$.  The normal form has degree at most $T$.

Start with a general polynomial $P$ of degree at most $L$.  Set the
residual in layer zero to $D_0=P$ and the other residuals to zero.
For $i=0,\ldots,n-1$, impose the linear condition
\[
 \operatorname{NF}_J(D_i)=0.
\]
Use \eqref{int:linear-division} for $D_i$.  Record
$a^i\sum I_\rho A_\rho(D_i)$ in the certificate.  The term
$a^{i+1}A(D_i)$ is added to the residual in layer $i+1$, or recorded
as a multiple of $a^n$ if $i+1=n$.  In the other term use the exact
identity
\[
 a^i f_0H(D_i)
  =a^i F_nH(D_i)
    -\sum_{j=1}^{n-1}a^{i+j}e_{m+j}(U)H(D_i).
\]
Record the first summand as a multiple of $F_n$.  For $i+j<n$ add
$-e_{m+j}H(D_i)$ to layer $i+j$; for $i+j\geq n$ record the term as
a multiple of $a^n$.  After the last layer no residual remains, so
these conditions imply the actual membership
\eqref{int:actual-membership}.  This procedure never takes a radical
or raises $P$ to a power.

All the residuals and conditions depend linearly on the coefficients
of $P$.  Choose $B$ so large that
\[
 \Delta+C_0(1+m+j)\leq Bj\quad(j\geq1),
 \qquad \Delta\leq B.
\]
Such $B$ is polynomial in $D$.  A contribution advancing by $j$
layers then increases residual degree by at most $Bj$; advancement
using the generator $a$ has the same bound with $j=1$.  Induction
gives
\begin{equation}\label{int:layer-degree}
 \deg D_i\leq L+Bi\quad(0\leq i<n).
\end{equation}
Every recorded term in the certificate also has degree polynomial
in $D+L+n$, including the terms advanced past the last layer.

To count the conditions, let $G$ be a polynomial bound for the degrees
of generators of the initial monomial ideal of $J$, and put
$r_0=\dim V(J)$.  In a surviving monomial, let $S$ be the set of
variables whose exponents are at least $G$.  No generator of the
initial ideal can be supported in $S$, since it would then divide the
monomial.  The coordinate space in the variables $S$ is therefore
contained in the zero set of that monomial ideal, so $|S|\leq r_0$.
There are at most $2^M$ choices of $S$, at most $G^M$ choices for
the remaining bounded exponents, and at most $(1+T)^{r_0}$ choices
of total degree at most $T$ in the unbounded exponents.  Hence the
space of normal forms of degree at most $T$ has dimension at most
\[
 2^MG^M(1+T)^{r_0}
   \leq B'(1+T)^{k-2}.
\]
The equality of the dimensions of an ideal and its initial ideal
follows from equality of their degree filtrations under a
degree-compatible order.  This count applies to nonradical ideals
as well.  If $J$ is the unit ideal its normal-form space is zero.
Combining the count with \eqref{int:layer-degree}, the entire
construction imposes at most the left side of
\eqref{int:condition-count} independent linear conditions.

On the other hand, the restrictions of degree-$\leq L$ polynomials
to $W'$ span a space of dimension at least $\binom{L+k}{k}$.
Indeed, $k$ generic linear coordinates are algebraically independent
on $W'$; their monomials of total degree at most $L$ are linearly
independent.  If every polynomial satisfying the layer conditions
vanished on $W'$, the codimension of the solution space in the
degree-$\leq L$ polynomial space would be at least this number.
That contradicts \eqref{int:condition-count}.  A polynomial nonzero
on $W'$ satisfying \eqref{int:actual-membership} therefore exists.

For the asserted height control, let $T$ be the polynomial certificate
degree bound just proved.  Introduce unknown coefficients for $P$
of degree at most $L$ and for the certificate polynomials of degree
at most $T$, and equate coefficients in
\[
 P=\sum_\rho B_\rho I_\rho+B_a a^n+B_f F_n.
\]
This is a homogeneous linear system of polynomial size in $D+L+n$.
The height of its coefficient matrix is polynomial in these quantities
times the stated input heights.  The minors argument in
\Cref{int:algebra} supplies a basis of solutions with the claimed
height bound.  At least one basis vector has its $P$ component nonzero
on $W'$; otherwise every solution would vanish there, contradicting
the existence already established.

Finally set $n=cL$.  For $k\geq2$ and $L\geq1$, the left side of
\eqref{int:condition-count} is at most
\[
 B'c(2+Bc+B)^{k-2}L^{k-1},
\]
whereas the right side is at least $L^k/k!$.  Thus it suffices to
choose the integer
\[
 L>k!B'c(2+Bc+B)^{k-2}.
\]
This is polynomial in $D+c$, because $k\leq M$ is fixed.  For $k=1$
there are no conditions and any $L\geq1$ suffices.  This proves all
assertions.
\end{proof}

\subsection{The product formula and dimension descent}

\begin{proof}[Proof of \Cref{int:height}]
All implicit constants in this proof depend on the fixed data in the
theorem.  We keep an irreducible closed $W\ni z$ whose defining
equation count and degrees are bounded by a polynomial in $d$ and
whose defining coefficient heights are bounded by
\begin{equation}\label{int:descent-control}
 D_i(d)(1+\log h)^{B_i},
\end{equation}
where $D_i$ is a fixed polynomial and $B_i$ is a fixed nonnegative
integer for the current descent step $i$.  Its degree has a polynomial
bound as well.  Initially $W=\mathbb A^N$ and these assertions hold.
At most $N$ strict dimension drops will be made.  Consequently
enlarging the polynomials and exponents at each step still yields
constants depending only on the initial fixed data.

First suppose $W\cap\{a=0\}$ is empty.  If $F_1,\ldots,F_r$ define
$W$ set theoretically, the Nullstellensatz says
$(F_1,\ldots,F_r,a)=(1)$.  This assertion is true for the displayed
ideal itself: an ideal with empty zero set is the unit ideal.
By \Cref{int:algebra} there is an identity
\[
 1=aA+\sum_{\rho=1}^r F_\rho A_\rho
\]
whose degrees are polynomial in $d$ and whose coefficient heights
have the form \eqref{int:descent-control}.  At $z$ we get $1=aA(z)$.
If $\tau_v$ is one at an archimedean place and zero otherwise, a
polynomial $Q$ of degree $T$ in $M$ variables satisfies
\begin{equation}\label{int:evaluation}
 \log|Q(x)|_v
  \leq \log\max(1,\|Q\|_v)
       +T\log\max(1,\|x\|_v)
       +\tau_v\log\binom{T+M}{M},
\end{equation}
where $\|Q\|_v$ denotes its maximum coefficient absolute value.
On $V$, $|a|_v<1$, so only the positive logarithms of $U$ enter this
bound.  Evaluation of $A$ and \eqref{int:local-input} give
\[
 \frac{h}{C_0d^{C_0}}
 \leq\sum_{v\in V}w_v\log|a|_v^{-1}
 =\sum_{v\in V}w_v\log|A(z)|_v
 \leq D_*(d)(1+\log h)^{B_*}.
\]
An inequality $h\leq D(d)(1+\log h)^B$ implies a polynomial bound
for $h$: for fixed $B$ there is a constant $c_B$ such that
$(1+\log h)^B\leq c_B\sqrt h$ for all $h\geq2$, and hence
$h\leq c_B^2D(d)^2$.  This proves the theorem in the present case.
It includes every case where $a$ is constant on $W$, since its value
is the nonzero target value, and every zero-dimensional $W$.

We may therefore suppose $a$ varies on $W$ and its boundary is
nonempty.  Condition \textup{(iv)} and \Cref{int:multiplicity} give
an integer combination $E$ and an integer $m$, with bounds polynomial
in $d$, satisfying \eqref{int:strict-order}.  Its coefficient and
tail bounds retain the forms required in the theorem.  More
explicitly, if $c_{\max}=\max_\alpha c_\alpha$, the ultrametric
inequality and integrality of the $c_\alpha$ preserve the original
tail bound at finite places.  At infinity replace $b_v$ by
\[
 b_v+\tau_v\log\max(1,s c_{\max}).
\]
The total extra weighted cost is at most
$\log\max(1,s c_{\max})$, since the total archimedean weight is one.
Combining the coefficient tuples likewise preserves their bounds.
We use $b_v$ for these enlarged errors.

Put $F=E_{<m}$ and form the graph variety $W'$ of
\Cref{int:graph}.  Its dimension is $k=\dim W$ and its equation
count, degrees, and heights have the bounds
\eqref{int:descent-control}, by \Cref{int:algebra}.  Lift the target
to $z'=(a,U,q)$, with $q=a^{-m}F(z)$.  The rational expression for
$q$ gives
\[
 h_{\mathrm{aff}}(z')
   \leq C_1(d)h+C_2(d)(1+\log h)^{B_2}
   \leq C_3(d)h,
\]
where $C_3$ is polynomial in $d$.  For the second inequality use
$(1+\log h)^{B_2}\leq c_{B_2}h$; its constant is fixed at this
descent step.  The polynomial $C_3$ is chosen before the interpolation
degrees $n,L$.

The local bound for $q$ is much smaller.  Since $E(z)=0$ on $V$,
\[
 q=-a^{-m}\sum_{j\geq m}a^je_j(U),
 \qquad
 \log^+|q|_v\leq(m+1)b_v.
\]
Consequently
\begin{equation}\label{int:lifted-local}
 \sum_{v\in V}w_v\log\max(1,\|z'\|_v)
    \leq C_4(d)(1+\log h)^{B_4}.
\end{equation}
The exact equation at the lifted target is
\[
 q+\sum_{j\geq0}a^je_{m+j}(U)=0.
\]

Take set-theoretic defining equations for $W'$ and let $I$ be the
ideal they generate.  The ideal is kept as given, without assuming
it radical.  By \Cref{int:graph},
\[
 J=(I,a,q+e_m(U))
 \quad\text{satisfies}\quad\dim V(J)\leq k-2.
\]
Set $A_0(d)=C_0d^{C_0}$, and first choose an integer
\begin{equation}\label{int:ratio-choice}
 c\geq 2A_0(d)(C_3(d)+1).
\end{equation}
Choose $L$ by \Cref{int:auxiliary} and set $n=cL$.
Both choices are polynomial in $d$, independent of $h$.  That lemma
gives $P$ of degree at most $L$, nonzero on $W'$, with a certificate
\begin{equation}\label{int:certificate}
 P=\sum_\rho B_\rho I_\rho+B_a a^n
       +B_f\left(q+\sum_{j=0}^{n-1}a^je_{m+j}(U)\right).
\end{equation}
Its degrees are polynomial in $d$.  The coefficient heights of $P$
and every $B_\rho,B_a,B_f$ are at most
$C_5(d)(1+\log h)^{B_5}$.  To apply the lemma with many coefficients,
bound the height of their joint tuple by the sum of the individual
heights in \eqref{int:coefficient-input}; there are polynomially many
indices $m,\ldots,m+n-1$, so this still has the asserted form.

We now test whether $P(z')$ can be nonzero.  At a selected place, the
parenthesized expression in \eqref{int:certificate} is
\[
 -a^{-m}\sum_{j\geq m+n}a^je_j(U).
\]
Its absolute value is at most
$|a|_v^n\exp((m+n+1)b_v)$.  Evaluate the two remaining certificate
terms by \eqref{int:evaluation}, use
\eqref{int:lifted-local}, and include $\log2$ at infinity for their
sum.  If $P(z')\ne0$, the result is
\begin{equation}\label{int:selected-smallness}
 \sum_{v\in V}w_v\log|P(z')|_v
 \leq -n\sum_{v\in V}w_v\log|a|_v^{-1}
          +C_6(d)(1+\log h)^{B_6}.
\end{equation}
Here polynomial certificate degrees multiply only the small positive
local logarithms, and $(m+n+1)\sum_Vw_vb_v$ is of the same permitted
size.  This explains precisely where
\eqref{int:local-input} is needed in addition to the global height
bound.

At all other places evaluate $P$ itself, of degree at most $L$,
using \eqref{int:evaluation}.  Its certificate is unnecessary for
this estimate.  Summing gives
\begin{equation}\label{int:outside-growth}
 \sum_{v\notin V}w_v\log|P(z')|_v
   \leq L C_3(d)h+C_7(d)(1+\log h)^{B_7}.
\end{equation}
Enlarge the number field to contain $W$, $W'$, their chosen
components, and all coefficients in \eqref{int:certificate}; the
normalized sums and all previously established estimates are
unchanged.  The product formula applied to the nonzero algebraic
number $P(z')$ combines \eqref{int:selected-smallness} and
\eqref{int:outside-growth} into
\[
 0\leq\left(-\frac n{A_0(d)}+LC_3(d)\right)h
                   +C_8(d)(1+\log h)^{B_8}.
\]
By \eqref{int:ratio-choice} the coefficient of $h$ is at most $-L$.
It follows that $h\leq C_8(d)(1+\log h)^{B_8}$ and hence that $h$
is polynomially bounded in $d$ as above.

If this polynomial height bound does not hold, we have proved
$P(z')=0$.  Substitute $q=F/a^m$ and clear denominators:
\[
 H(a,U)=a^{mL}P(a,U,F(a,U)/a^m).
\]
This is a polynomial.  It vanishes at $z$ because $a(z)\ne0$, and
it is not identically zero on $W$, because $P$ is nonzero on the
irreducible graph closure and its graph over $a\ne0$ is dense.
Its degree is at most $mL+L\max(1,\deg F)$; expanding the
substitution bounds its coefficient height by
$h_{\mathrm{aff}}(P)+Lh_{\mathrm{aff}}(F)$ plus a polynomial degree
error.  Thus it has exactly the controls
\eqref{int:descent-control}.  Choose an irreducible component
containing $z$ of the proper intersection $W\cap\{H=0\}$.
By \Cref{int:algebra} its defining equations have the same types of
degree and height bounds, and its dimension is strictly smaller.

Repeat with this component as $W$.  Condition \textup{(iv)} applies
to it because it is an irreducible algebraic variety containing the
same target.  After at most $N$ drops one either has obtained the
polynomial height bound or reaches the empty-boundary case, which
gives that bound.  There are finitely many fixed-step polynomial
bounds in this procedure, so their maximum is bounded by $Cd^C$
for a constant depending only on the fixed data.  This completes
the proof.
\end{proof}
